% 指针生成网络在解码第t步的计算：编码器状态、注意力分布、上下文向量、生成门与混合输出分布。
% 注意力数值为教学示例，与正文算例一致。
\begin{tikzpicture}[
  x=1mm,y=1mm,figure picture,
  cell/.style={figure node,minimum width=13mm,minimum height=6mm,
    inner sep=1.5pt,font=\figurefont\scriptsize},
  tok/.style={font=\figurefont\scriptsize},
  flow/.style={figure flow},
  note/.style={font=\figurefont\scriptsize,text=FigureMuted},
  row label/.style={font=\figurefont\scriptsize,text=FigureMuted,anchor=east},
  rnn/.style={draw=FigureStroke,line width=.6pt,
    {Stealth[length=1.1mm]}-{Stealth[length=1.1mm]}},
  bar/.style={fill=FigureInputFill,draw=FigureInput,line width=.5pt},
  halo/.style={preaction={draw=white,line width=2.6pt}}
]
  % 源端词元、编码器状态与注意力分布
  \foreach \i/\w/\a in {1/读书会/0.05,2/改至/0.10,3/海棠楼/0.25,4/201/0.60}{
    \pgfmathsetmacro{\xx}{(\i-1)*18}
    \node[tok] (x\i) at (\xx,0) {\w};
    \node[cell,draw=FigureInput,fill=FigureInputFill] (h\i) at (\xx,10) {$\bm h_{\i}$};
    \draw[flow] (x\i) -- (h\i);
    \pgfmathsetmacro{\hh}{\a*20}
    \draw[bar] (\xx-3,20) rectangle (\xx+3,20+\hh);
    \node[note,anchor=south,inner sep=1pt] at (\xx,20+\hh) {\a};
  }
  \draw[rnn] (h1) -- (h2);
  \draw[rnn] (h2) -- (h3);
  \draw[rnn] (h3) -- (h4);
  \draw[figure link] (-4,20) -- (58,20);
  \node[row label] at (-6,0) {源词元};
  \node[row label] at (-6,10) {编码器状态};
  \node[row label] at (-6,25) {注意力$\bm a_{t}$};

  % 上下文向量
  \node[cell,draw=FigureModel,minimum width=34mm] (ctx) at (27,46)
    {上下文$\bm h^{*}_{t}=\sum_i a_{t,i}\bm h_i$};
  \draw[flow] (27,36) -- (ctx.south);

  % 解码器
  \node[tok] (yprev) at (88,0) {上一词元“海棠楼”};
  \node[cell,draw=FigureModel] (s) at (88,10) {$\bm s_t$};
  \draw[flow] (yprev) -- (s);
  \draw[figure link,densely dashed,-{Stealth[length=1.2mm]}] (s.west)
    -- node[note,above,sloped,pos=.45] {查询} (61,23);

  % 词表分布
  \node[cell,draw=FigureModel,minimum width=28mm] (pv) at (88,46)
    {词表分布$P_{\mathrm{vocab}}$};
  \draw[flow] (s.north) -- (pv.south);
  \draw[flow] (ctx.east) -- (pv.west);

  % 生成门
  \node[circle,draw=FigureLoss,fill=FigureLossFill,line width=.75pt,
    minimum size=7mm,inner sep=0pt,font=\figurefont\scriptsize] (g) at (58,62) {$\sigma$};
  \node[note,anchor=west] at (g.east) {生成门$p_{\mathrm{gen}}$};
  \draw[flow] (ctx.north east) -- (g);
  \draw[flow,halo] (s.north west) -- (g.south east);

  % 混合输出分布
  \node[cell,draw=FigureOutput,fill=FigureOutputFill,minimum width=72mm] (out) at (58,80)
    {$p(w)=p_{\mathrm{gen}}P_{\mathrm{vocab}}(w)+(1-p_{\mathrm{gen}})\sum_{i:x_i=w}a_{t,i}$};
  \draw[flow] (g.north) -- (g.north |- out.south);
  \draw[flow] (pv.north) -- node[note,right] {生成} (pv.north |- out.south);
  \draw[flow] (-4,22) -- node[note,left,pos=.55] {复制} (-4,80) -- (out.west);
\end{tikzpicture}
